\section{Introduction}
\label{sec:intro}

%% [§1 ¶1]
Large language models (LLMs) can hallucinate, rely on outdated knowledge, and lack
specialised facts. Retrieval-augmented generation (RAG)~\cite{lewis2020rag} provides
the generator with passages from an external corpus. Combining dense search with a
cross-encoder re-ranker further improves question-answering (QA) accuracy across many
LLMs~\cite{yang2025enhancing}.

%% [§1 ¶2]
Multi-hop questions combine facts from multiple documents, for example to identify a
bridge entity and ask about its
properties~\cite{yang2018hotpotqa,ho2020twowiki,trivedi2022musique}. A single retrieval
often misses second-hop evidence. Agents therefore decompose questions into sub-questions
and gather evidence iteratively. Small language models (SLMs) running locally often struggle to
control this loop. They may answer too early, retrieve after collecting enough evidence,
or ask irrelevant sub-questions.

%% [§1 ¶3]
These control decisions are small, structured prediction problems: \emph{Is the evidence
sufficient? What should happen next? How much evidence is still missing?} Using the
generator ties decision quality to model size and consumes tokens at every
step. Laya~\cite{laya} is a non-autoregressive \dq{System~1} decision model. It answers
multiple-choice, yes/no, and ordinal questions in one forward pass, returning a
probability for each option. Training uses strictly proper scoring rules, and temperature
scaling calibrates these probabilities. Most adaptive retrieval methods use the
generator's signals or decide once per query (Section~\ref{sec:related}). When external
calibrated control should replace the generator's control, and how this depends on
generator size, remain open questions.

%% [§1 ¶4]
Our $2\times2$ factorial design uses a shared decomposition loop to isolate the
controller's effect. We vary the \emph{evidence source}
(closed-book recall or RAG with dense retrieval and re-ranking) and the \emph{loop
controller} (the LLM or Laya). Laya's decisions are enforced at every hop, rather than
offered as an optional tool. Thinking mode is disabled for all models. We evaluate five
Qwen3.5 sizes~\cite{qwen35} on three multi-hop benchmarks, with 200 questions per
configuration ($5\times4\times3=60$ configurations). We report exact match (EM),
token-level F1, cost, paired bootstrap 95\% confidence intervals (CIs), and the
controller's Brier score and expected calibration error (ECE). We ask
three research questions (RQs).
\begin{itemize}
  \item \textbf{RQ1 (Accuracy).} With the decomposition loop fixed, does replacing the
        generator's control with a calibrated external decision model change multi-hop QA
        accuracy (EM and F1)? Does the effect depend on the evidence source?
  \item \textbf{RQ2 (Cost).} How does external control affect tokens, hops, end-to-end
        latency, and tool-call reliability per question?
  \item \textbf{RQ3 (Scale).} How do these effects change with generator scale, from
        0.8B to 27B parameters?
\end{itemize}

%% [§1 ¶5]
Our four contributions address these research questions as follows.
\begin{itemize}
  \item A factorial framework separating evidence and controller effects within one loop,
        with wrong answers classified at the tool, retrieval, stopping, or answering
        stage (RQ1).
  \item A weak-supervision method that simulates agent states from supporting-paragraph
        annotations to train a calibrated multi-hop controller for RQ1--RQ3.
  \item An analysis of tokens, hops, latency, and tool-call reliability under explicit
        control in local SLM agents (RQ2).
  \item An evaluation of five Qwen3.5 sizes on three datasets showing that the accuracy
        effect of external control reverses with generator scale (RQ3).
\end{itemize}

%% [§1 ¶6]
Sections~\ref{sec:related}--\ref{sec:experiments} cover related work, methods, and results.
Section~\ref{sec:discussion} answers the RQs and discusses limitations.
Section~\ref{sec:conclusion} concludes the paper.
